% Introduction (Gate A; deepened with literature at the user's request after Gate C): five paragraphs, about 650 words;
% NARRATIVE_V5_CELL.md, Introduction plan. Four open questions, each set against the debate it belongs to, then "Here, we" with
% the findings in figure order (no figure call-outs here: Figure 1 is first called out in the Results).

Complete connectomes now map every synapse of an insect's brain and nerve cord\cellcite{dorkenwald2024,berg2026,bates2026}, and
their cell types can be matched across animals and between the sexes\cellcite{schlegel2024,sturner2025}. Models built on them
predict sensorimotor pathways and neural activity, either without training or with training confined to one sensory
pathway\cellcite{shiu2024,lappalainen2024,pospisil2024}. In a physics model of the fly's body, trained controllers walk and fly,
but they are not built from the fly's wiring\cellcite{vaxenburg2025}. The connectome itself is often described as a static map
without dynamics, state or plasticity\cellcite{kyllmann2026}. Whether a fixed wiring diagram can learn demanding behavior, and
what its memories would be made of, is unknown.

How a brain holds information across a delay is still disputed. Persistent firing carries remembered items in recordings from
monkeys and humans\cellcite{thrower2023,paluch2025}, in areas from prefrontal to visual cortex\cellcite{curtis2021}, whereas
activity-silent states\cellcite{stokes2015} can be read out by an impulse\cellcite{wolff2017}. Whether silent states hold working
memories, serve only passive storage or reflect retrieval from episodic memory is
debated\cellcite{trubutschek2019,beukers2021,oberauer2022}. Such states are almost always attributed to synapses, through
short-term facilitation or fast Hebbian plasticity\cellcite{mongillo2008,masse2019,pals2020,lansner2023}, while intrinsic states
such as spike-frequency adaptation have been studied mainly as aids to computation\cellcite{salaj2021}. In the fly, visual
memories that last several seconds depend on signaling in ring neurons of the central complex\cellcite{neuser2008,kuntz2017}.
Which substrate learning recruits on real wiring is open.

Where memories live, and why learned codes hold in some states but not in others, is equally unresolved. Engrams of a single
memory are spread across many brain regions\cellcite{josselyn2020,roy2022}, and the excitability of a neuron at learning biases
whether it joins one\cellcite{delamare2024,hansel2024}. Population codes can stay stable across a delay while responses to an
impulse change with time\cellcite{murray2017,wolff2020}; stable and dynamic codes coexist in prefrontal
cortex\cellcite{sapountzis2022}, and in human visual cortex the code for a remembered location turns toward the coming eye
movement\cellcite{li2023}. Retrieval can require the brain state present at encoding\cellcite{jovasevic2015,meyer2017}. Changes of
excitability alone can break a decoder trained in one state, whereas a decoder trained across states keeps
working\cellcite{haimerl2025}. Which part of a memory a change of state alters, the stored trace or its reading, has not been
tested in a complete nervous system.

In Alzheimer's disease, early memory loss can be a failure of access rather than of storage. In mouse models of early disease,
activating engram cells recalls memories that natural cues no longer retrieve\cellcite{roy2016}, and activating parvalbumin
interneurons restores retrieval\cellcite{park2024}. Neuronal hyperexcitability is an early feature of the
disease\cellcite{zott2019,palop2016,li2024,okechukwu2026}, but lowering it has had mixed effects on cognition in
patients\cellcite{bakker2012,vossel2021,faini2025}. In flies, amyloid-β42 (Aβ42) degrades the A-type potassium channel
Kv4\cellcite{ping2015}, and the loss of Kv4 from mushroom-body neurons underlies part of the courtship memory deficit that Aβ42
causes\cellcite{feng2018}. Why a lesion that spares new learning removes existing memories is not understood. Each of these
questions could be asked of a connectome that learns, in which every state variable of every cell can be recorded and moved
between brain states.

Here, we trained a digital twin of the fly's nervous system: the complete wiring of one male animal, simulated cell by cell.
It was trained on tasks rather than fitted to recorded activity, checked against fly physiology and replayed bit for bit, with
every experiment declared before it ran. The model learns to plan moves and to remember a visual cue. Spikes write the cue into
neuronal adaptation, which holds it silently until the next input reads it out. The trace is spread across the nervous system,
and no ensemble of cells is more necessary than matched sets. The cells that hold it are enriched for transcripts of
\gene{SK}, a calcium-activated potassium channel, in fly cell atlases. The learned readout uses the trace only at the delays it
learned, because the same trace drives a different response at each delay. The wiring keeps the readout reachable under heavy
cell loss, and training across excitability states widens access, as decoder theory predicts. A disease-like loss of potassium current
after learning rewrites the trace and changes how it is read, so recall returns only in the state of learning, while learning
under the loss is spared.
